\section{为什么行业流出、公司却可能被提前布局}

徐工机械所在的机械设备行业本周净流出162.71亿元，但徐工自身连续5日净流入2.60亿元，同期仅涨0.12\%。这种逆行业资金行为必须由基本面解释，否则可能只是被动护盘。2026Q1公司营收297.91亿元，同比增长9.26\%；归母净利润20.56亿元，同比增长0.86\%，看似利润增速弱，但经营活动现金流净额21.52亿元，同比增长153.12\%。利润受人民币升值和财务费用扰动，现金回收和收入扩张却在改善。

\section{增长来源}

\begin{exhibitbox}[Exhibit 7：徐工机械增长与质量桥]
\begin{tabularx}{\textwidth}{L{3.4cm}R{2.4cm}X}
\toprule
\textbf{指标/驱动} & \textbf{数据} & \textbf{含义} \\
\midrule
2026Q1收入 & 297.91亿元/+9.26\% & 主业延续增长 \\
2026Q1归母净利 & 20.56亿元/+0.86\% & 汇率和财务费用压制表面增速 \\
经营现金流 & 21.52亿元/+153.12\% & 应收和支出管理显著改善 \\
2025海外收入 & 486亿元/+16.6\% & 占收入约48.2\%，国际化进入主导阶段 \\
2025矿业机械收入 & 93.8亿元/+47.4\% & 第二增长曲线 \\
股权激励考核 & 2026净利不低于75亿元或两年累计不低于140亿元 & 管理层目标提供底线，但不是业绩保证 \\
\bottomrule
\end{tabularx}
\sourcenote{公司2026Q1报告；华创、东吴、太平洋证券2026年业绩点评。}
\end{exhibitbox}

真正需要验证的是：汇率扰动是否暂时、海外高增长是否能抵消国内周期、现金流改善是否来自可持续的经营质量而非阶段性回款。公司2025年融资租赁回购义务显著压降，有助于降低风险敞口；但海外占比接近一半，也使贸易壁垒和汇率成为更重要变量。

\section{价格、资金与拥挤度}

7月10日收盘8.43元，位于过去一年价格区间约20\%，距离区间高点回撤约32\%。融资余额约8.42亿元，占流通市值约1.12\%，处于可控水平；7月9日融资略有净偿还，说明连续主力流入并非融资盘全面加杠杆。北向在2026Q1前十大流通股东中明显增持，但二季度公开口径显示持仓市值存在波动，因此报告将其定义为机构参与而非单向锁仓。

\section{估值与动作}

三份原始券商报告给出的2026E EPS区间约0.67--0.72元。AStock取0.71元，基准16倍PE对应11.36元；叠加9.40元市场锚和华创12.50元外部目标，最终目标11.08元，相对8.43元空间约31.4\%。

\begin{itemize}
  \item \textbf{配置区}：8.18--8.50元，分批而非一次性重仓。
  \item \textbf{确认区}：站稳9.20--9.60元，且中报海外收入和现金流继续改善。
  \item \textbf{失效条件}：2026E EPS下修至0.60元以下、海外增速低于8\%、经营现金流重新显著转弱。
\end{itemize}
